% Generated draft: requires mathematical and source review.
\documentclass[11pt]{article}
\usepackage[margin=25mm]{geometry}
\usepackage{fontspec}
\setmainfont{Latin Modern Roman}
\usepackage{amsmath,amssymb,mathtools,unicode-math}
\setmathfont{Latin Modern Math}
\usepackage{enumitem,needspace,xurl,hyperref}
\setlength{\parindent}{0pt}
\setlength{\parskip}{5pt}
\setlength{\emergencystretch}{3em}
\newcommand{\N}{\mathbb{N}}
\newcommand{\Z}{\mathbb{Z}}
\newcommand{\Q}{\mathbb{Q}}
\newcommand{\R}{\mathbb{R}}
\newcommand{\C}{\mathbb{C}}
\newcommand{\F}{\mathbb{F}}
\newcommand{\A}{\mathbb{A}}
\newcommand{\PP}{\mathbb{P}}
\newcommand{\E}{\mathbb{E}}
\newcommand{\eps}{\varepsilon}
\newcommand{\dd}{\,\mathrm{d}}
\begin{document}
\section*{UnsolvedMath ID: 700809}
\textbf{Source problem code:} catalogue entry
\subsection{Definitions and mathematical statement}
Use the standard finite-group convention for Fitting classes. A finite group \(H\) is soluble if its derived series, defined by \(H^{(0)}=H\) and \(H^{(i+1)}=[H^{(i)},H^{(i)}]\), reaches the trivial group. A nonempty Fitting class \(\mathfrak X\) is a nonempty class of finite groups, closed under isomorphisms, taking normal subgroups, and taking products of normal subgroups: if \(A,B\trianglelefteq H\) and \(A,B\in\mathfrak X\), then \(AB\in\mathfrak X\). It is called soluble if every member is soluble.

For a finite group \(H\), a subgroup \(U\le H\) is \(\mathfrak X\)-maximal in \(H\) if \(U\in\mathfrak X\) and every subgroup \(W\) satisfying \(U\le W\le H\) and \(W\in\mathfrak X\) equals \(U\). A subgroup \(N\le G\) is subnormal if there is a finite chain \(N=N_0\trianglelefteq N_1\trianglelefteq\cdots\trianglelefteq N_r=G\), allowing \(r=0\). An \(\mathfrak X\)-injector of \(G\) is a subgroup \(V\le G\) such that \(V\cap N\) is \(\mathfrak X\)-maximal in \(N\) for every subnormal subgroup \(N\) of \(G\).

\textbf{Question (L. A. Shemetkov, 1990).} For every nonempty soluble Fitting class \(\mathfrak X\) and every finite non-soluble group \(G\), does there exist an \(\mathfrak X\)-injector \(V\le G\)?

The supplied record does not define these technical terms; the definitions above make their standard finite-group interpretation explicit.
\subsection{Short English statement}
Does every finite non-soluble group possess an injector for each nonempty Fitting class consisting of soluble finite groups? Such an injector must intersect every subnormal subgroup in a subgroup maximal among those belonging to the class. The supplied registry labels this question open; that label is not independently verified here.
\subsection{Sources}
\begin{itemize}
\item Supplied local UnsolvedMath registry record, permanent ID \(700809\), exact source problem code \texttt{catalogue entry}, titled \emph{Kourovka Notebook Problem 11.117}. The record attributes the problem to L. A. Shemetkov (1990) and labels it open.
\item \emph{The Kourovka Notebook (Unsolved Problems in Group Theory)}, original problem 11.117, 11th issue (1990); the supplied background identifies the 21st-issue edition (2026), \texttt{arXiv:1401.0300v48}. The supplied collection description names E. I. Khukhro and V. D. Mazurov as editors. These bibliographic details are supplied attribution and have not been independently verified.
\end{itemize}
\end{document}
